\documentclass[11pt]{article}

\usepackage[T1]{fontenc}
\usepackage{lmodern}
\usepackage[margin=0.90in]{geometry}
\usepackage{microtype}
\usepackage{amsmath,amssymb,amsthm,mathtools,mathrsfs}
\usepackage{booktabs,longtable,array,tabularx}
\usepackage{graphicx}
\usepackage{float}
\usepackage{xcolor}
\usepackage{enumitem}
\usepackage{fancyhdr}
\usepackage{hyperref}
\usepackage[nameinlink,noabbrev]{cleveref}
\usepackage{listings}
\usepackage{caption}
\usepackage{bm}
\usepackage{needspace}
\usepackage{placeins}
\usepackage{siunitx}

\graphicspath{{figures/}}
\allowdisplaybreaks
\sisetup{scientific-notation=true,round-mode=places,round-precision=4}

\hypersetup{
  pdftitle={Projector-Resolved Seed Calculus X},
  pdfauthor={Komorebi P.},
  pdfsubject={Universal residual seeds, deformation forests, directed transport, seed zeta functions, companion completion, compatible carriers, and finite occupancy},
  colorlinks=true,
  linkcolor=black,
  citecolor=black,
  urlcolor=black
}

\pagestyle{fancy}
\fancyhf{}
\fancyhead[L]{Projector-Resolved Seed Calculus X}
\fancyhead[R]{Komorebi P.}
\fancyfoot[C]{\thepage}
\setlength{\headheight}{14pt}

\newtheorem{theorem}{Theorem}[section]
\newtheorem{proposition}[theorem]{Proposition}
\newtheorem{lemma}[theorem]{Lemma}
\newtheorem{corollary}[theorem]{Corollary}
\theoremstyle{definition}
\newtheorem{definition}[theorem]{Definition}
\newtheorem{example}[theorem]{Example}
\newtheorem{conjecture}[theorem]{Conjecture}
\newtheorem{classicalinput}[theorem]{Classical Input}
\newtheorem{numericalcertificate}[theorem]{Numerical Certificate}
\theoremstyle{remark}
\newtheorem{remark}[theorem]{Remark}

\newcommand{\Ocal}{\mathcal O}
\newcommand{\mfrak}{\mathfrak m}
\newcommand{\F}{\mathbf F}
\newcommand{\Z}{\mathbf Z}
\newcommand{\Q}{\mathbf Q}
\newcommand{\C}{\mathbf C}
\newcommand{\Spf}{\operatorname{Spf}}
\newcommand{\Spec}{\operatorname{Spec}}
\newcommand{\GL}{\operatorname{GL}}
\newcommand{\Hom}{\operatorname{Hom}}
\newcommand{\End}{\operatorname{End}}
\newcommand{\Tr}{\operatorname{Tr}}
\newcommand{\Frob}{\operatorname{Frob}}
\newcommand{\WD}{\operatorname{WD}}
\newcommand{\HT}{\operatorname{HT}}
\newcommand{\ad}{\operatorname{ad}}
\newcommand{\red}{\operatorname{red}}
\newcommand{\Def}{\operatorname{Def}}
\newcommand{\rad}{\operatorname{rad}}
\newcommand{\srank}{\operatorname{srank}}
\newcommand{\Cl}{\operatorname{Cl}}
\newcommand{\SCC}{\operatorname{SCC}}
\newcommand{\rank}{\operatorname{rank}}
\newcommand{\im}{\operatorname{im}}
\newcommand{\coker}{\operatorname{coker}}
\newcommand{\id}{\operatorname{id}}
\newcommand{\Sys}{\mathscr S}
\newcommand{\Forest}{\mathfrak F}
\newcommand{\Compat}{\mathfrak C}
\newcommand{\Packet}{\mathfrak P}

\setlength{\parindent}{0pt}
\setlength{\parskip}{0.62em}
\setlist[itemize]{leftmargin=1.55em,itemsep=0.20em,topsep=0.30em}
\setlist[enumerate]{leftmargin=1.75em,itemsep=0.24em,topsep=0.32em}
\captionsetup{font=small,labelfont=bf}

\lstset{
  basicstyle=\ttfamily\scriptsize,
  breaklines=true,
  columns=fullflexible,
  frame=single,
  showstringspaces=false,
  tabsize=4,
  numbers=left,
  numberstyle=\tiny,
  xleftmargin=1.5em,
  framexleftmargin=1.2em
}

\title{\textbf{Projector-Resolved Seed Calculus X}\\[0.35em]
\large Universal Residual Seeds: Deformation Forests, Transport Categories, and System Completion}
\author{Komorebi P.}
\date{}

\begin{document}
\maketitle

\begin{abstract}
A residual representation is a genuine zeroth-order coordinate.  All positive-depth perturbations above it may fill its formal deformation neighborhood, but they cannot move into a different residual class.  This paper develops the corresponding universal object: a residual forest whose roots are residual representations, whose finite levels are lifts modulo powers of a uniformizer, and whose infinite branches are complete characteristic-zero objects.  We prove the residual partition, no-crossing, inverse-limit reconstruction, finite-branch compactness, exact smooth and singular branch counts, packet decomposition of tangent and obstruction cohomology, a directed minimum-seed theorem, rational seed-zeta formulas, and finite companion-forest completion.  The coefficient prime two is represented by a weighted forest in which characteristic-zero packets may merge modulo two and become visible only at later depth.  We then define a prime-independent system seed through a common Frobenius spine, prove a compatible-carrier summand theorem under finite Frobenius separation, and reduce finite occupancy in a fixed semisimple carrier to the nonnegative integer trace of an exact joint projector.  The paper does not prove that every admissible residual companion exists or that every finite Hecke packet is occupied.  Those two existence statements remain the exact horizontal boundary beyond vertical deformation theory.
\end{abstract}

\begin{center}
\fbox{\begin{minipage}{0.94\linewidth}
\textbf{Block 19 formal-verification status.} Fifteen source proof claims remain checked and five remain partial at full arithmetic scope. A sufficiently deep complete orthogonal family of equivariant quotient projectors now lifts to an actual decomposition inside a supplied finite-group carrier, with precision loss independent of the number of labels. The actual invariant images and their reductions are checked. Over local coefficients with a nonunit ideal generator, lifts of the same positive-depth family are equivariantly conjugate by a unit congruent to the identity. The intended arithmetic packet and period-data interpretation remains unfinished.
\end{minipage}}
\end{center}


\paragraph{Eighth-block formalization status.}
The cumulative formal ledger records 15 source proof claims checked at their
stated scope, 5 partially formalized, and no untouched proof claims. Partial
claims 8.2, 9.1, 9.2, 11.1, and 11.2 are not completed arithmetic theorems.
This status does not prove any universal existence conjecture. No PDF has been
typeset or visually certified for this code-focused checkpoint.

\tableofcontents
\newpage

\section{The universal seed at residual depth}

The earlier papers treat a seed as a primitive object together with an operator grammar, a packet system, a topology, and an exact reconstruction law.  In arithmetic deformation problems, the primitive coordinate has two layers:

\[
\boxed{
\text{residual root}
\quad+\quad
\text{positive-depth deformation coordinates}.
}
\]

The distinction is structural.  A correction divisible by the uniformizer changes higher digits but leaves the residual digit fixed.  Therefore the correct universal arithmetic seed cannot be one residual representation unless the admissible residual space has only one point.  It must contain every residual root and the complete formal neighborhood above each root.

Let \(\Ocal\) be a complete discrete valuation ring with uniformizer \(\varpi\), maximal ideal \((\varpi)\), residue field
\[
k=\Ocal/(\varpi),\qquad |k|=q,
\]
and fraction field \(K\).  When a characteristic-zero conclusion is intended, assume additionally that \(\operatorname{char}K=0\); the general counting statements do not require that restriction.  Let \(\mathcal M\) be a moduli functor on complete local Noetherian \(\Ocal\)-algebras with residue field \(k\), equipped with a well-defined reduction map
\[
\red_A:\mathcal M(A)\longrightarrow \mathcal M(k).
\]
For a Galois representation this means either that a framed integral model has been fixed or that the residual semisimplified class is used.

\begin{definition}[Residual seed system]
For \(\bar x\in\mathcal M(k)\), define the residual fiber
\[
\mathcal M_{\bar x}(A)
=
\{x\in\mathcal M(A):\red_A(x)=\bar x\}.
\]
The point \(\bar x\) is the residual root.  The family of all nonempty fibers is the residual special fiber, and the formal deformation problems above them form the residual forest.
\end{definition}

For a Galois problem the principal example is
\[
\bar\rho:G_{F,S}\longrightarrow \GL_n(k),
\]
with a deformation functor \(\Def_{\bar\rho}\) cut out by a determinant and selected local conditions.

\begin{figure}[H]
\centering
\includegraphics[width=0.87\linewidth]{01_residual_seed_architecture.pdf}
\caption{The universal residual seed is not one point.  It is the special fiber together with one complete formal neighborhood above each admissible root.}
\end{figure}

\subsection{What the paper is designed to separate}

Three statements are often compressed into one sentence even though they have different logical content.

\begin{enumerate}
\item \textbf{Vertical completeness:} every compatible perturbation stream above a fixed root reconstructs a point of its complete deformation space.
\item \textbf{Horizontal companion existence:} every required residual root, including roots at other coefficient primes, actually exists.
\item \textbf{Occupancy:} the resulting arithmetic packet occurs in a nonzero global carrier of the required type.
\end{enumerate}

The first belongs to deformation and inverse-limit theory.  The second is a residual companion problem.  The third is an arithmetic or automorphic nonvanishing problem.  Neither of the last two follows from arbitrarily deep perturbation above one already occupied root.

\section{Residual partition and no crossing}

\begin{theorem}[Residual partition]
For every coefficient algebra \(A\),
\[
\boxed{
\mathcal M(A)
=
\bigsqcup_{\bar x\in\mathcal M(k)}
\mathcal M_{\bar x}(A).
}
\]
\end{theorem}

\begin{proof}
Every \(A\)-valued object has one reduced object \(\red_A(x)\in\mathcal M(k)\), so the union is exhaustive.  Two distinct fibers cannot intersect because an object cannot have two different reductions.  Hence the union is disjoint.
\end{proof}

\begin{theorem}[No crossing under positive-depth perturbation]
Suppose \(x,y\in\mathcal M(\Ocal)\) are related by a deformation operation that is congruent to the identity modulo \(\varpi\).  Then
\[
\red(x)=\red(y).
\]
Every finite composition remains in one residual fiber.  If an infinite perturbation stream converges in the \(\varpi\)-adic topology, its limit remains in the same fiber.  The existence of that limit is an explicit hypothesis, not a consequence of positive depth alone.
\end{theorem}

\begin{proof}
A correction in \(\varpi\Ocal\) vanishes after reduction modulo \(\varpi\).  Thus one operation preserves the residual root, and induction proves the finite statement.  For a convergent stream, convergence at the first finite precision makes its reduction eventually equal to that of its limit.  Every term already has the initial reduction, so the limit has that reduction too.  Completeness applies to Cauchy streams; it does not make every root-preserving stream Cauchy.  For example, \(0,3,0,3,\ldots\) has constant reduction modulo \(3\) but no limit modulo \(9\).
\end{proof}

\begin{corollary}[Zeroth-order obstruction]
If two admissible objects have different residual roots, no sequence of genuinely positive-depth perturbations beginning at one can reach the other.  A horizontal transport between their roots is a zeroth-order operation, not a deeper lift.
\end{corollary}

\begin{figure}[H]
\centering
\includegraphics[width=0.95\linewidth]{03_no_crossing_branches.pdf}
\caption{Each perturbation stream remains above its initial residual root.  Refinement adds digits vertically; it does not move a branch horizontally.}
\end{figure}

\begin{numericalcertificate}[Exact residual partition]
The finite implementation partitions the integers \(0,\ldots,8\) into the three classes modulo \(3\):
\[
\{0,3,6\},\qquad \{1,4,7\},\qquad \{2,5,8\}.
\]
Every tested correction of the form \(x\mapsto x+p h\) preserves the residual class exactly.
\end{numericalcertificate}

\section{Lift forests and inverse-limit reconstruction}

Fix \(\bar x\in\mathcal M(k)\).  For \(m\ge1\), set
\[
X_m(\bar x)=\mathcal M_{\bar x}(\Ocal/\varpi^m).
\]
Reduction gives transition maps
\[
r_{m+1,m}:X_{m+1}(\bar x)\longrightarrow X_m(\bar x).
\]
A compatible branch is a sequence \((x_m)_{m\ge1}\) satisfying
\[
r_{m+1,m}(x_{m+1})=x_m.
\]

\begin{classicalinput}[Represented complete deformation problems]
When a deformation functor is represented by a complete local ring \(R_{\bar x}\), its \(\Ocal\)-valued points are continuous local maps \(R_{\bar x}\to\Ocal\).  Completeness identifies such maps with compatible maps to every quotient \(\Ocal/\varpi^m\).
\end{classicalinput}

\begin{theorem}[Formal inverse-limit reconstruction]
Assume \(\mathcal M_{\bar x}\) is represented by a complete local \(\Ocal\)-algebra \(R_{\bar x}\).  Then reduction induces a natural bijection
\[
\boxed{
\mathcal M_{\bar x}(\Ocal)
\cong
\varprojlim_m
\mathcal M_{\bar x}(\Ocal/\varpi^m).
}
\]
\end{theorem}

\begin{proof}
A continuous local homomorphism \(R_{\bar x}\to\Ocal\) reduces to a compatible family of homomorphisms into \(\Ocal/\varpi^m\).  Conversely, compatible finite-level maps determine a map
\[
R_{\bar x}\longrightarrow\varprojlim_m\Ocal/\varpi^m\cong\Ocal.
\]
The two constructions are inverse because equality in a separated complete ring is detected modulo every power of its maximal ideal.
\end{proof}

\begin{theorem}[Finite-branch compactness]
Suppose every \(X_m(\bar x)\) is finite and nonempty and every node has finitely many children.  Then the lift tree above \(\bar x\) contains an infinite compatible branch.
\end{theorem}

\begin{proof}
The tree is infinite because it has a nonempty level at every depth, and it is finitely branching by hypothesis.  K\"onig's lemma gives an infinite path.  The edge relation is reduction, so that path is a compatible branch.
\end{proof}

Surjectivity of every transition map is sufficient but not necessary.  Finite branching and nonempty arbitrarily deep levels are the compactness mechanism.

\begin{proposition}[Smooth formal branch count]
If
\[
R_{\bar x}\cong\Ocal[[t_1,\ldots,t_d]]\quad\text{as complete local }\Ocal\text{-algebras},
\]
then
\[
\boxed{
\#X_m(\bar x)=q^{d(m-1)}.
}
\]
Fix a uniformizer and a section \(s:k\to\Ocal\) of reduction, with \(s(0)=0\). Every \(\Ocal\)-valued branch is encoded by \(d\) positive-depth digit streams
\[
t_i=s(a_{i,1})\varpi+s(a_{i,2})\varpi^2+\cdots,
\qquad a_{i,j}\in k,
\]
\end{proposition}

\begin{proof}
A local map sends each parameter into the image of \((\varpi)\) in \(\Ocal/\varpi^m\). Conversely, every tuple in that image is nilpotent and defines a coefficient-preserving evaluation homomorphism. This homomorphism is local. To prove uniqueness without assuming continuity of an arbitrary competing homomorphism, observe in one variable that a homomorphism killing \(X^N\) only sees the finite polynomial truncation modulo \(X^N\). Induction over the finite number of variables gives uniqueness in all dimensions. The actual local maps are therefore in bijection with the full parameter tuple set. Its cardinality is \(q^{d(m-1)}\), including the zero-dimensional chart.

For integral maps, reduce an arbitrary tuple in \((\varpi)\Ocal\) at every precision, evaluate there, and use the represented-map reconstruction theorem. The finite evaluations are compatible, the reconstructed algebra map is local, and every competing local map has the same finite reductions. Separatedness gives uniqueness. Local maps are continuous in the maximal-ideal adic topologies. This establishes the integral parameter classification rather than assuming it from the finite count.

For the digit claim, subtract the chosen representative of the residue and divide the remainder by the uniformizer. Iterating gives unique digits. Conversely, the partial sums of any digit stream form a compatible Cauchy family and reconstruct uniquely by completeness and separatedness. More precisely, for each parameter there is a unique stream \((a_{i,j})_{j\ge1}\) such that
\[
\varpi^{n+1}\mid t_i-\sum_{j=1}^{n}s(a_{i,j})\varpi^j\qquad(n\ge0).
\]
The section is a set-theoretic section of reduction, not an asserted embedding of residue fields into the coefficient ring. Finally, a supplied local \(\Ocal\)-algebra isomorphism from \(R_{\bar x}\) transports both the map classification and the count.
\end{proof}

\begin{figure}[H]
\centering
\includegraphics[width=0.94\linewidth]{02_residual_digit_forest.pdf}
\caption{The first four levels of a full \(3\)-adic digit forest.  A path is a compatible finite-depth perturbation stream.}
\end{figure}

\begin{figure}[H]
\centering
\includegraphics[width=0.82\linewidth]{04_inverse_limit_digits.pdf}
\caption{A compatible base-five residue tower reconstructs the digit stream \((2,4,1,3,0,2)\).}
\end{figure}

\section{Tangent and obstruction packets}

For a residual representation \(\bar\rho\), first-order deformations and lifting obstructions are classically controlled by cohomology of the adjoint module.  The projector calculus does not replace that theory.  It resolves the controlling complex into invariant packets when suitable projectors exist.

\begin{classicalinput}[Tangent and obstruction spaces]
For a framed representation problem, first-order lifts are described by \(Z^1(G,\ad\bar\rho)\).  For an unframed Schur problem, tangent directions are described by \(H^1(G,\ad\bar\rho)\).  Relations in a presentation of a deformation ring map into \(H^2(G,\ad\bar\rho)\), which contains the obstruction classes.
\end{classicalinput}

Let
\[
C^0\xrightarrow{d_0}C^1\xrightarrow{d_1}C^2\xrightarrow{}\cdots
\]
be a cochain complex over a field.  Suppose a finite family of projectors \(P_j^r:C^r\to C^r\) satisfies
\[
\sum_jP_j^r=I,
\qquad
P_i^rP_j^r=\delta_{ij}P_j^r,
\qquad
d_rP_j^r=P_j^{r+1}d_r.
\]

\begin{theorem}[Cohomology packet decomposition]
Under these hypotheses, each packet is a subcomplex and
\[
\boxed{
H^r(C^\bullet)
\cong
\bigoplus_jH^r(P_jC^\bullet).
}
\]
In particular,
\[
H^1=\bigoplus_jH^1_j,
\qquad
H^2=\bigoplus_jH^2_j,
\]
so tangent and obstruction spaces reconstruct exactly from their packet components.
\end{theorem}

\begin{proof}
For each packet, use the actual image submodules
\(C_j^r=\operatorname{im}P_j^r\). Commutation with the differential
makes these a subcomplex. Write
\[
Z^r=\ker d_r,\qquad B^r=\operatorname{im}d_{r-1}\subseteq Z^r,
\qquad H^r=Z^r/B^r,
\]
and define \(Z_j^r,B_j^r,H_j^r\) using the restricted differentials of
\(C_j^\bullet\). The inclusion \(i_j:C_j^\bullet\to C^\bullet\)
and the image-valued projection \(p_j:C^\bullet\to C_j^\bullet\)
are cochain maps, so they induce maps on these genuine quotients.

For \(i\ne j\), orthogonality gives
\(H(p_i)H(i_j)=0\). Idempotence gives
\(H(p_j)H(i_j)=I\), and completeness gives
\(\sum_j H(i_j)H(p_j)=I\). Thus the maps
\[
[z]\longmapsto([P_j^r z])_j,
\qquad
([z_j])_j\longmapsto\left[\sum_j i_j z_j\right]
\]
are inverse linear maps. The finite product equals the direct sum, proving
the assertion. In the formalization, both sides are also identified with
Mathlib's standard homology objects of the genuine graded complexes.
\end{proof}

\paragraph{Scope of the splitting.}
The algebraic proof works over any commutative ring and does not divide by
\(2\), so it applies in characteristic two when the stated integral
chain projectors exist. It does not construct those projectors. The
identification of a particular Galois tangent or obstruction problem with
this cohomology remains the preceding classical input.
The splitting detects vanishing:
\[
H^r(C^\bullet)=0\quad\Longleftrightarrow\quad
H^r(C_j^\bullet)=0\text{ for every }j.
\]
It does not force either side to vanish or decouple nonlinear deformation
equations. Moreover, commutation with the differential is essential:
complementary idempotents on the underlying vector spaces need not
preserve cocycles or coboundaries.

\paragraph{Naturality.}
If a cochain map \(F:C^\bullet\to D^\bullet\) intertwines the supplied
projectors, it restricts to maps \(F_j:C_j^\bullet\to D_j^\bullet\).
The direct-sum decomposition commutes with \(H^r(F)\) and the family
\(H^r(F_j)\). This is checked using representatives in the actual
quotients, independently of any dimension computation.


\begin{remark}[Odd and dyadic nonlinear equations]
At odd residue characteristic, a differential graded Lie model may be written in Maurer--Cartan form
\[
d\eta+\frac12[\eta,\eta]=0,
\]
and a cyclic grading converts this into exact packet convolutions.  At residue characteristic two, division by \(2\) is not integral.  A quadratic or divided-square refinement must be retained.  This is one reason the dyadic residual forest is weighted rather than obtained by formal substitution into the odd-prime formulas.
\end{remark}

\begin{numericalcertificate}[Finite packet complex]
The exact modular test uses two invariant subcomplexes over \(\F_5\).  Their cohomology dimensions are
\[
(0,0,0),
\qquad
(1,2,1),
\]
and the reconstructed total is
\[
\dim(H^0,H^1,H^2)=(1,2,1).
\]
The packet-sum residual is zero.
\end{numericalcertificate}

\paragraph{New checked realization of the finite pattern.}
The Block 04 implementation separately uses
\[
\mathbf F_5^2\xrightarrow{d_0}\mathbf F_5^4\xrightarrow{d_1}\mathbf F_5^2,
\quad d_0(a,b)=(a,0,0,0),\quad d_1(a,b,c,d)=(b,0).
\]
The first packet is the exact complex with dimensions \((1,2,1)\),
and the second packet has zero differentials and dimensions \((1,2,1)\).
The actual middle cohomology is linearly isomorphic to \(\mathbf F_5^2\),
which is checked in Lean. Independent exhaustive coset tests check all
three cohomology degrees after changes of basis. This is a new specified
realization of the displayed dimension pattern, not a claimed rerun of
the original companion implementation.


\begin{figure}[H]
\centering
\includegraphics[width=0.78\linewidth]{11_tangent_obstruction_packets.pdf}
\caption{A finite chain-projector model.  Tangent and obstruction dimensions are carried by one packet and reconstruct exactly.}
\end{figure}

\section{The weighted dyadic residual forest}

At an odd prime, distinct sign characters remain distinct in the residue field.  At the coefficient prime two,
\[
1\equiv-1\pmod2.
\]
Thus distinct characteristic-zero packets can share one residual root and separate only after more precision is retained.

\begin{definition}[Visibility weight]
For two integral packet coordinates \(a,b\in\Ocal\), define their separation depth
\[
w(a,b)=\min\{m\ge1:a\not\equiv b\pmod{\varpi^m}\},
\]
when such an \(m\) exists.  A dyadic packet decomposition is weighted by these first visibility depths and by the valuations of any denominators required by its rational projectors.
\end{definition}

\begin{example}[The two sign branches]
The characters \(\chi_+(g)=1\) and \(\chi_-(g)=-1\) have the same reduction modulo \(2\), but
\[
1\not\equiv-1\pmod4.
\]
Their first separation depth is therefore \(2\).  One mod-two root contains two branches that become distinct at the next level.
\end{example}

\begin{figure}[H]
\centering
\includegraphics[width=0.82\linewidth]{05_dyadic_delayed_separation.pdf}
\caption{The sign packets merge at residual depth and separate at depth two.  This is a genuine change in integral coordinates, not a numerical plotting effect.}
\end{figure}

A simple weighted formal model already captures delayed visibility.  Let \(w_i\ge1\) and constrain
\[
t_i\in(\varpi^{w_i}\Ocal+\varpi^m\Ocal)/\varpi^m\Ocal.
\]
Then coordinate \(i\) is invisible until \(m>w_i\).

\begin{proposition}[Weighted smooth-chart count]
For weights \(\bm w=(w_1,\ldots,w_d)\),
\[
\boxed{
N_{\bm w}(m)
=
q^{\sum_{i=1}^d\max(m-w_i,0)}.
}
\]
\end{proposition}

\begin{proof}
The image of \(\varpi^{w_i}\Ocal\) in \(\Ocal/\varpi^m\) has one element when \(m\le w_i\), and \(q^{m-w_i}\) elements when \(m>w_i\).  Multiplying the independent coordinate counts gives the formula.
\end{proof}

This model does not assert that every dyadic Galois deformation ring is a weighted formal ball.  It gives the exact local counting law once such weighted coordinates or a Smith-normal-form chart have been proved.

\begin{figure}[H]
\centering
\includegraphics[width=0.82\linewidth]{07_weighted_formal_chart.pdf}
\caption{Weights delay finite-depth visibility.  Once all coordinates are active, the asymptotic slope equals the ordinary formal dimension.}
\end{figure}

\section{Seed zeta functions and branching dimensions}

Finite-depth counts record more than the existence of a branch.  They distinguish smooth, weighted, and singular local geometry.

\begin{definition}[Seed zeta function]
For a residual root \(\bar x\) with finite levels, define
\[
\boxed{
Z_{\bar x}(T)
=
\sum_{m\ge1}\#X_m(\bar x)T^m.
}
\]
When a weighted chart is specified, write \(Z_{\bm w}(T)\).
\end{definition}

\begin{theorem}[Smooth seed zeta function]
For a smooth formal component of dimension \(d\),
\[
\boxed{
Z_d(T)=\frac{T}{1-q^dT}.
}
\]
\end{theorem}

\begin{proof}
By the proved formal-model map-to-parameter equivalence, \(\#X_m=q^{d(m-1)}\) counts all the actual local coefficient-algebra maps. Summing the geometric series gives the identity. The checked statement is an equality of infinite formal power series, with zero constant coefficient.
\end{proof}

\begin{theorem}[Weighted seed zeta function]
Let \(1\le w_1\le\cdots\le w_d\), put \(W=\sum_iw_i\), and let
\[
a_m=q^{\sum_i\max(m-w_i,0)}.
\]
Then
\[
\boxed{
Z_{\bm w}(T)
=
\sum_{m=1}^{w_d}a_mT^m
+
q^{-W}\frac{(q^dT)^{w_d+1}}{1-q^dT}.
}
\]
In particular, every weighted smooth chart has a rational seed zeta function.
\end{theorem}

\begin{proof}
The finite sum records the depths before all coordinates are active.  For \(m\ge w_d+1\),
\[
a_m=q^{dm-W},
\]
so the remaining terms form one geometric series.
\end{proof}

\begin{theorem}[Singular power seed]
For integers \(e\ge1\) and \(m\ge1\), in \(\Ocal/\varpi^m\) the equation
\[
x^e=0
\]
has exactly
\[
\boxed{
q^{m-\lceil m/e\rceil}
}
\]
solutions.  Its seed zeta function is
\[
\boxed{
Z_e(T)
=
\frac{\sum_{r=1}^{e}q^{r-1}T^r}
{1-q^{e-1}T^e}.
}
\]
\end{theorem}

\begin{proof}
The congruence is equivalent to
\[
e\,v_\varpi(x)\ge m,
\]
so \(x\) must lie in \(\varpi^{\lceil m/e\rceil}\Ocal/\varpi^m\), which has the stated cardinality.  Write \(m=ek+r\), \(1\le r\le e\).  The exponent becomes \((e-1)k+r-1\).  Summing over \(k\ge0\) and \(1\le r\le e\) gives the rational expression.
\end{proof}

\begin{definition}[Branching dimension]
When the limit exists, define the finite-level counting exponent
\[
\dim_{\mathrm{seed}}(\bar x)
=
\lim_{m\to\infty}
\frac{\log_q\#X_m(\bar x)}{m}.
\]
A limsup may be used in general.
\end{definition}

\begin{corollary}[Fractional singular dimension]
For \(x^e=0\),
\[
\boxed{
\dim_{\mathrm{seed}}=1-\frac1e.
}
\]
In particular, \(x^2=0\) has branching dimension \(1/2\).

\paragraph{Counting exponent versus compatible branches.}
The preceding exponent counts all finite-level solutions, including nodes which cannot lift indefinitely. It is not the cardinality or dimension of the space of infinite branches. In fact every compatible tower of solutions to \(x^e=0\), for \(e\ge1\), is the zero tower. For each \(m\), a solution at depth \(em\) lies in \(\varpi^m\Ocal/\varpi^{em}\); its reduction at depth \(m\) is therefore zero. Compatibility proves the assertion at every depth. This observation and the real logarithmic limit are both formalized in Block 02.

\end{corollary}

\begin{figure}[H]
\centering
\includegraphics[width=0.83\linewidth]{06_smooth_singular_counts.pdf}
\caption{A smooth two-dimensional formal chart and the singular equation \(x^2=0\) have visibly different exponential growth.}
\end{figure}

\begin{figure}[H]
\centering
\includegraphics[width=0.78\linewidth]{08_fractional_seed_dimension.pdf}
\caption{The finite-depth branching dimension of \(x^e=0\) is \(1-1/e\), giving a family of fractional local growth exponents.}
\end{figure}

\section{Residual transport categories}

Vertical lifting never changes a residual root.  Horizontal coverage requires operations whose residual effect is not congruent to the identity.  Examples that may define such transports include twisting, duality, restriction, induction, tensor product, symmetric powers, exterior powers, and functorial transfer.  The word ``may'' is essential: each operation must preserve the deformation conditions and the realization statement being transported.

\begin{definition}[Admissible residual transport]
An admissible transport from \(\bar\rho\) to \(\bar\rho'\) consists of
\[
\mathcal T:\bar\rho\longmapsto\bar\rho'
\]
together with natural maps
\[
\mathcal T_A:\Def_{\bar\rho}(A)\longrightarrow\Def_{\bar\rho'}(A)
\]
for every coefficient algebra \(A\), compatible with base change and reduction.  When the transport is reversible, it is an isomorphism in the residual transport groupoid.  One-way operations give a directed category.
\end{definition}

For a finite collection of roots, record a directed edge \(u\to v\) whenever occupancy or the relevant property is known to propagate from \(u\) to \(v\).

\begin{definition}[Directed seed rank]
For a directed transport graph \(G\), define
\[
\srank_G
=
\min\{|S|:\Cl_G(S)=V(G)\},
\]
where \(\Cl_G(S)\) is the set of vertices reachable from \(S\).
\end{definition}

\begin{theorem}[Directed minimum-seed theorem]
Let \(G\) be a finite directed graph.  Collapse every strongly connected component to one vertex, obtaining the condensation DAG.  Then
\[
\boxed{
\srank_G
=
\#\{\text{source strongly connected components of }G\}.
}
\]
One representative from each source component is a minimum seed set.
\end{theorem}

\begin{proof}
Every vertex of a finite DAG is reachable from a source, so one seed in each source component reaches the entire condensation graph.  Within a strongly connected component, one vertex reaches every other.  Hence the source representatives are sufficient.

Conversely, a source component has no incoming edge from another component.  It cannot be reached from a seed outside itself, so every generating seed set must meet it.  Therefore the number of source components is also a lower bound.
\end{proof}

For reversible transports the condensation has one vertex per connected component, recovering the undirected minimum-seed theorem.

\paragraph{Formal scope and propagation.}
Reachability is the reflexive-transitive closure of the supplied directed edge
relation, and a strongly connected component is the actual quotient by mutual
reachability. The minimum seed count is defined independently as the least
cardinality of a set that reaches all vertices. Every covering set must meet
every source component; conversely, meeting every source component suffices
in a finite graph. Thus any choice of one representative from each source is
optimal. Empty graphs have rank zero, and self-loops are permitted. Adding
admissible edges cannot increase the minimum seed count.

If a property $P$ propagates along every supplied edge and holds at all selected
seeds, it holds at every reached vertex. These are explicit hypotheses. The
minimum-seed theorem does not construct valid arithmetic edges or supply
initial occupancy.

\paragraph{Why the finite hypothesis cannot be dropped.}
Take vertices $\{*\}\sqcup\{v_n:n\in\mathbf N\}$, with $*$ isolated and
edges $v_{n+1}\to v_n$. Every strongly connected component is a singleton,
and $\{*\}$ is the unique source component. No finite seed set covers all
vertices: it can reach only indices no larger than the largest chosen index.
Both the unique-source statement and the failure of every finite cover are
checked in Lean. Hence one source component in an infinite condensation does
not imply that one seed reaches all objects.


\begin{figure}[H]
\centering
\includegraphics[width=0.91\linewidth]{09_directed_transport_graph.pdf}
\caption{A finite directed transport model.  The components \(\{A,B\}\), \(\{D,E\}\), and \(\{G\}\) are sources; \(C\) and \(F\) are reached downstream.}
\label{fig:transport-graph}
\end{figure}

\begin{figure}[H]
\centering
\includegraphics[width=0.86\linewidth]{10_source_scc_condensation.pdf}
\caption{The condensation DAG makes the minimum seed count visible without searching every subset of roots.}
\end{figure}

\begin{numericalcertificate}[Directed seed count]
The test graph has five strongly connected components and three source components.  The theorem gives seed rank three.  Exhaustive search over every subset of vertices independently returns three, and the selected representatives reach all seven roots.

\paragraph{New explicit checked realization.}
Separately from the unretrieved original companion implementation, Block 05
specifies vertices $A,B,C,D,E,F,G$ and the eight edges
\[
 A\to B,\ B\to A,\ B\to C,\ C\to F,\ D\to E,\ E\to D,\ E\to F,\ G\to F.
\]
The source components are $\{A,B\}$, $\{D,E\}$, and $\{G\}$; the other
components are $\{C\}$ and $\{F\}$. A bounded path search to length three
is proved closed under every original edge, which certifies equality with
unbounded reachability. The resulting label certificate is identified with
the actual quotient by mutual reachability. Lean checks seed rank three and
exactly the four minimum sets
\[
 \{A,D,G\},\quad\{A,E,G\},\quad\{B,D,G\},\quad\{B,E,G\}.
\]
Independent exact tests enumerate every directed graph on zero through four
vertices, including loops, and every seed subset. These finite tests audit the
implementation; the universal theorem is a separate formal proof.

\end{numericalcertificate}

\section{Prime-independent system seeds}

A compatible system should not be treated as an unrelated collection of \(\lambda\)-adic roots.  Its members share algebraic Frobenius polynomials, Hodge--Tate data, determinant information, and local Weil--Deligne packets.

\begin{definition}[Common arithmetic spine]
Let \(M\) be a number field.  A common spine \(\Compat\) consists of:
\begin{itemize}
\item characteristic polynomials \(P_v(X)\in M[X]\) at almost all places;
\item coefficient-prime independent weight and determinant data;
\item prescribed local Weil--Deligne or inertial packet labels at the ramified places.
\end{itemize}
For each finite place \(\lambda\) of \(M\), let
\[
X_{\lambda,m}(\Compat)
\]
be the finite set of admissible objects modulo \(\varpi_\lambda^m\) realizing the reduction of this spine.
\end{definition}

The system seed is the fibered object
\[
\boxed{
\Sys_{\mathrm{sys}}
=
\Sys_2^{\mathrm{wt}}
\times_{\Compat}
\prod_{\lambda\nmid2}'\Sys_\lambda.
}
\]
The fiber product notation states a compatibility requirement.  It does not prove that every fiber is nonempty.

\begin{theorem}[Adelic finite-depth completion]
Assume:
\begin{enumerate}
\item an algebraic common spine \(\Compat\) over a number field \(M\) is fixed;
\item for every finite place \(\lambda\) of \(M\) and every \(m\ge1\), the companion level \(X_{\lambda,m}(\Compat)\) is finite and nonempty;
\item the reduction forests are finitely branching;
\item the selected local conditions are closed under inverse limits.
\end{enumerate}
Then for every \(\lambda\) there is an infinite compatible branch
\[
\rho_\lambda\in\varprojlim_mX_{\lambda,m}(\Compat),
\]
and the collection \(\{\rho_\lambda\}_\lambda\) realizes the prescribed strictly compatible spine.
\end{theorem}

\begin{proof}
Finite-branch compactness gives an infinite branch in every coefficient-prime forest.  Inverse-limit closure gives an admissible \(\lambda\)-adic object.  By definition each finite level has the prescribed reductions of the common Frobenius and local packets.  Equality at every finite precision implies equality in the complete coefficient ring.  Thus all members realize the same spine.
\end{proof}

\begin{corollary}[Residual companion obstruction]
If
\[
X_{\lambda,1}(\Compat)=\varnothing
\]
for one coefficient prime \(\lambda\), no compatible completion with spine \(\Compat\) exists.  Deeper perturbations at other primes cannot create the missing residual companion.
\end{corollary}

\begin{figure}[H]
\centering
\includegraphics[width=0.82\linewidth]{13_common_system_spine.pdf}
\caption{The local members are coordinates of one system only after their finite data descend to a common algebraic spine.}
\end{figure}

\begin{example}[A scalar companion failure]
For
\[
f(x)=x^2+1,
\]
there are two simple branches at \(p=5\) and \(p=13\), no residual roots at \(p=3\) or \(p=7\), and one root modulo \(2\) that has no child modulo \(4\).  A complete \(5\)-adic branch therefore does not imply an adelic solution.
\end{example}

\begin{figure}[H]
\centering
\includegraphics[width=0.79\linewidth]{12_companion_forest_occupancy.pdf}
\caption{Finite companion levels for \(x^2+1=0\).  The obstruction at \(2,3,7\) is visible before any deep lifting argument begins.}
\end{figure}


\paragraph{Checked representation-level completion and its arithmetic boundary.}
The eighth formalization block works with actual matrix representations of a
fixed group, not only with underlying sets. For a complete separated ring $R$
and ideal $I$, compatible maps
$\rho_m:G\to\operatorname{GL}_d(R/I^m)$ reconstruct a unique integral
representation for that fixed compatible family. Multiplication and inverses
are verified after every reduction. If all finite representations are continuous
for the discrete quotient topology, the integral representation is continuous
for the ideal-adic topology. This does not assert uniqueness among all possible
branches of the allowed forest.

The forest has as its level the actual subtype of representations satisfying
the declared finite conditions. Those conditions must be preserved by reduction.
Finiteness and nonemptiness at every level give an infinite branch and hence an
actual continuous representation. Finiteness of the coefficient quotient alone
is not a substitute for finiteness of the allowed representation set for an
arbitrary source group. Surjectivity of the reductions is not required.

All coefficients of the characteristic polynomial commute with reduction.
Consequently, prescribed integral characteristic polynomials at every precision
hold exactly in the reconstructed representation. For a common coefficient field
$E$, local integral polynomials lie in $\Ocal_\lambda[X]$ and are compared with
$P_v\in E[X]$ only after mapping both into the local coefficient field. No field
embedding $E\to\Ocal_\lambda$ is assumed. The allowed-place predicate is explicit
and may depend on $\lambda$. Trivial action on a supplied subgroup is likewise
detected at every precision; identifying that subgroup with arithmetic inertia
is a separate input.

The eleventh block assembles continuous coefficient-field representations, exact
common polynomials and the prescribed local predicate from positive finite
depths. Depth-zero finiteness and nonemptiness are derived. Claims 8.2 and 11.1
already assume inverse-limit closure of the selected local conditions; an
unconditional closure theorem is not a missing hypothesis of their conditional
proofs. Their conditional matrix core is checked. The remaining full arithmetic
interpretation must identify the admissible finite objects and the actual local
invariants with these data. The coherent-framing obligation is resolved by Block 14
below. An arbitrary local
predicate is not being renamed Hodge--Tate or Weil--Deligne compatibility.


\paragraph{Completion from conjugacy classes (Block 14).}
Over an $I$-adically complete separated commutative coefficient ring $R$,
for a fixed finite rank, the formalization now takes finite-level objects to be
actual matrix representations modulo conjugation by an invertible matrix.
Assume the coefficient quotients $R/I^k$ are finite.
Conjugation commutes with coefficient reduction, and each conjugacy class has
finitely many frames over a finite coefficient ring. Thus finite allowed class
levels give finite full preimages of allowed matrix representations. Downward
admissibility and positive-depth nonemptiness yield an exactly compatible
matrix branch and an actual integral representation. More strongly, every
prescribed compatible class branch has such representatives, using finite
nonempty fibers over that branch. No surjectivity of reductions or lifting of
prescribed change-of-basis matrices is assumed. This does not assert uniqueness
up to integral conjugacy.

For finitely generated $I$ with finite $R/I$, finiteness of $R/I^k$ is derived.
One continuous representative per allowed class suffices, because continuity
is proved invariant under conjugation. Characteristic polynomials, determinants
and trivial action on a supplied subgroup are likewise actual invariants.
The unframed common-spine theorem constructs continuous coefficient-field
representations with the prescribed whole polynomials and the explicitly
supplied inverse-limit-closed local predicate.

An actual $\mathbb Z_3$ example applies the class-level completion theorem
for the cyclic group of order two and polynomial $X^2-1$, producing a nonidentity generator. Another
checked example has compatible classes represented by
$\operatorname{diag}(1,-1)$ modulo $3$ and
$\operatorname{diag}(-1,1)$ modulo $9$: these chosen matrices disagree under
reduction, and an actual swap of basis repairs them. Arbitrary independent
representatives therefore do not replace the framing argument.

The framing obligation is finished. Identifying the manuscript's arithmetic
admissible objects, actual local invariants and finite-level packet reductions
with this interface remains separate. An uninterpreted local predicate is not
a construction of Hodge--Tate or Weil--Deligne compatibility. The source's local
closure hypothesis remains explicit, and Claim 8.2 remains partial.


\paragraph{Actual local comparisons from fixed-depth certificates (Block 15).}
Let $R$ be an $I$-adically complete separated commutative ring with finite
coefficient quotients and an injective map into a field $K$. Fix a finite matrix
rank. For finitely many
local group maps $\phi_s:H_s\to G$ and integral template representations
$\tau_s$, fix detecting depths $d_s$. At precision $m$ retain matrices $P_{s,m}$
as part of the admissible object and require
\[
\rho_m(\phi_s(h))P_{s,m}=P_{s,m}\tau_{s,m}(h).
\]
For $m\ge d_s$ require their determinants to remain nonzero after reduction to
$R/I^{d_s}$. Earlier levels omit that test. Finite occupied augmented levels
give a compatible branch of both representations and comparisons. The integral
limits satisfy the actual equations and have nonzero determinants, so the
comparisons are units over $K$. A converse supplies all finite nodes from any
integral realization with these detectors. Finite-level invertibility and
compatible prior choices of comparisons are not assumed.

The unframed positive-depth version also retains continuity and whole common
polynomials. Its actual local isomorphisms are derived from the certificates;
there is no uninterpreted local-stability predicate for this particular problem.
The fixed depths and augmented-level occupancy remain explicit inputs.

The formalization specializes to the actual maps
$G_{\mathbb Q_p}\to G_{\mathbb Q}$ induced by the field embeddings and choices
of algebraic closures. A nontrivial continuous quadratic character of
$G_{\mathbb Q}$, constructed from its action on $i^2=-1$, gives the integral
$2$-adic actions
\[
\tau(\sigma)=\operatorname{diag}(1,-1),\quad
\rho(\sigma)=\begin{pmatrix}1&-1\\0&-1\end{pmatrix},\quad
P=\begin{pmatrix}1&1\\0&2\end{pmatrix}
\]
at a nonidentity quadratic element. The determinant $2$ is detected modulo $4$.
The representations are conjugate over $\mathbb Q_2$ but not integrally
conjugate. A separate $\mathbb Z_3$ example invokes simultaneous completion
and reconstructs a comparison which is an integral nonunit and a field unit.

The bound is substantive: for the unipotent actions
$\rho_a(g)=\begin{pmatrix}1&ag\\0&1\end{pmatrix}$, every integral intertwiner
with $\rho_1$ has determinant $aP_{22}^2$. Taking $a=\pi^n$ for a nonzero
$\pi\in I$ rules out every
fixed detecting depth, even though all the field-valued approximants are
isomorphic and their integral limit is trivial. This does not refute Claim 8.2,
which already assumes inverse-limit closure of its chosen local conditions.

Exact prescribed local Galois action is stronger than an inertial-type label.
The common arithmetic packet interpretation, its coefficient comparisons and
the actual Hodge--Tate/Weil--Deligne attachment remain separate. The new theorem
does not supply a uniform bound for all intended arithmetic packets. Claim 8.2
remains partial.

\paragraph{Derived comparison bounds (Block 16).}
For a finite group $\Gamma$ whose order is a unit in a commutative local
coefficient ring $R$, any prescribed equivariant residual isomorphism lifts to
an integral isomorphism. Lift its matrix to $A$ and define
\[
P=|\Gamma|^{-1}\sum_{g\in\Gamma}\rho(g)A\tau(g^{-1}).
\]
Reindexing proves equivariance, and reduction recovers the prescribed matrix.
Its determinant is a unit. Complete local coefficients with finite quotients,
occupied compatible finite conditions and this residual type therefore give
actual integral completion with detector depth one. A fixed continuous map
$G\to\Gamma$ permits pullback to infinite source groups. The finite-level
representations must factor through the same map, and the residual type is
imposed on their $\Gamma$-actions. No uniform determinant bound or augmented
certificate occupancy is assumed in this sector.

For a second sector let $R$ be a DVR with uniformizer $\pi$, fraction field $K$
and matrix rank $n$. If a fixed template $\tau$ has full $K$-linear action span,
finite generation and denominator clearing give a fixed $c$ such that
\[
\pi^c M_n(R)\subseteq\operatorname{span}_R\{\tau(g):g\in G\}.
\]
Normalize a field comparison to an integral $P$ with a unit entry. The displayed
span constructs an integral $B$ with $BP=\pi^cI$, so
\[
\det P\mid\pi^{cn},\qquad \det P\not\equiv0\pmod{\pi^{cn+1}}.
\]
One $c$ works for all integral representations in the fixed field-isomorphism
class. Jacobson density derives full span from actual irreducibility and a
scalar commutant. The formalization states these hypotheses directly and does
not assume the DVR fraction field is algebraically closed.

If at every precision there is an integral representation of this rational
type satisfying the finite conditions, the derived bound supplies the
simultaneous certificates of Block 15 and preserves the type after completion.
The integral approximants are not required to be coherent, but their existence
is a substantive premise. Arbitrary torsion nodes are not asserted to lift.

Over $\mathbb Z_3$, the actual average corrects
$A=\left(\begin{smallmatrix}1&1\\0&1\end{smallmatrix}\right)$ to
$P=\left(\begin{smallmatrix}1&-1/2\\0&1\end{smallmatrix}\right)$ for the
involutions $\rho(g)=\left(\begin{smallmatrix}1&1\\0&-1\end{smallmatrix}\right)$
and $\tau(g)=\operatorname{diag}(1,-1)$. It preserves the residual comparison
and has determinant one. Pullback through the quadratic character of the
actual $G_{\mathbb Q}$ gives continuous arithmetic examples. In contrast, the
identity and $\operatorname{diag}(1,\mathrm{sign})$ over $\mathbb Z/4$ agree
modulo two but are not integrally conjugate. Their group order is not a unit.
These restricted results do not supply general arithmetic packet attachment,
all finite-level existence statements, or cross-coefficient compatibility.

\section{Compatible carriers and finite Frobenius separation}

One useful completion mechanism is available when the desired packet already occurs as a semisimple summand of a known compatible carrier.

\begin{theorem}[Finite Frobenius separation]
Let $\mathscr W_1,\ldots,\mathscr W_r$ be pairwise nonisomorphic strictly
compatible systems whose members are absolutely irreducible, with common
algebraic unramified trace values $t_i(v)$ in a number field $E$.
There are $r$ distinct permitted unramified places $v_1,\ldots,v_r$ for which
\[
 C=(t_i(v_j))_{j,i}\in M_r(E),\qquad \det C\ne0.
\]
Every coefficient-field embedding $\sigma_\lambda:E\hookrightarrow E_\lambda$
therefore sends $C$ to an invertible matrix. Its entries agree with the
ordinary representation traces
$\operatorname{Tr}(W_{i,\lambda}(\operatorname{Frob}_{v_j}))$
when the selected places satisfy the unramified compatibility convention
for that coefficient prime. In particular, places above the residue
characteristic of $\lambda$ are excluded from this ordinary trace comparison.
\end{theorem}

\begin{proof}
Fix one coefficient prime. The classical independence of the characters of
pairwise nonisomorphic absolutely irreducible representations and the density
of unramified Frobenius conjugacy classes give independent restricted trace
functions. Density here concerns the union of the conjugacy classes, not an
unsupported assertion about a set of chosen representatives.

For $r$ independent functions, their evaluation linear forms span the
$r$-dimensional coefficient dual: a proper span would leave a nonzero
coefficient vector annihilated by every evaluation, contradicting independence.
Choose a basis from the evaluations. Its $r$ distinct sample points give an
invertible square matrix. The common algebraic traces put that matrix over
$E$; an injective field embedding reflects the nonvanishing of its determinant.
All further field embeddings preserve the same nonvanishing.
\end{proof}

\paragraph{Formalization boundary.}
Characteristic-zero independence is now Lean-checked for actual finite-dimensional
irreducible representations of arbitrary monoids. Finiteness of the group and
algebraic closedness of the field are not needed for this result. The proof
uses the actual monoid algebra and Jacobson density to select constituent
actions. Jointly continuous coordinate actions have continuous characters.
Coverage in every open normal quotient of a compact totally disconnected group
implies density of the conjugacy-saturated samples and hence the finite separator.
Arithmetic Frobenius coverage remains unformalized. The source proof also needs
the bridge from pairwise nonisomorphic systems to pairwise nonisomorphic members
at the fixed coefficient prime under the intended compatibility definition.
The full arithmetic theorem remains partial. Selection is existential and gives
no effective norm bound. A common algebraic matrix is still distinct from an
ordinary unramified trace comparison at all coefficient primes simultaneously.

\begin{theorem}[Compatible-carrier summand completion]
Suppose a semisimple strictly compatible carrier decomposes as
\[
\mathscr W
\cong
\bigoplus_{i=1}^{r}
\mathscr W_i^{\oplus c_i},
\]
with the \(\mathscr W_i\) as above.  At one coefficient prime \(\lambda_0\), suppose a semisimple summand has the form
\[
V_{\lambda_0}
\cong
\bigoplus_{i=1}^{r}
W_{i,\lambda_0}^{\oplus m_i},
\qquad
0\le m_i\le c_i.
\]
Then
\[
\boxed{
V_\lambda
=
\bigoplus_{i=1}^{r}
W_{i,\lambda}^{\oplus m_i}
}
\]
defines a strictly compatible summand at every coefficient prime.  The vector \((m_i)\) is uniquely recoverable from finitely many Frobenius traces supplied by the separation theorem.
\end{theorem}

\begin{proof}
Use the same multiplicities in every member of the compatible carrier.  Direct sums preserve strict compatibility, and \(m_i\le c_i\) makes \(V_\lambda\) a summand of \(W_\lambda\).  At the separating places, the target trace vector is \(C\bm m\).  Full column rank makes \(\bm m\) the unique solution.
\end{proof}

\paragraph{Recovery and realization are different steps.}
For a certified square separator the coefficient vector is $C^{-1}t$.
This is a valid carrier multiplicity vector only when every entry is a
nonnegative integer and satisfies its capacity bound $m_i\le c_i$.
Those tests and uniqueness of natural multiplicities in characteristic zero
are separately Lean-checked, as is transport through a field embedding.
Reduction modulo a residue characteristic is not a field embedding and may
lose rank. The representation construction and common-polynomial part are now
Lean-checked as described below. The full strictly compatible arithmetic
statement, including its local conditions, remains partially formalized.

\paragraph{Scope of the supplied-carrier hypothesis.}
Claim 9.2 supplies the compatible carrier, its decomposition and the compatible
constituent systems. Constructing unknown such objects is not required by that
conditional theorem. The remaining formal obligation is that finite sums preserve
the intended actual local arithmetic structures, together with the arithmetic
separator input. The algebraic monodromy results below do not by themselves
identify attached Weil--Deligne representations or Hodge--Tate comparison data.

\paragraph{Explicit summands inside a supplied decomposition.}
Let $K$ be a field, let $G$ be a monoid, and let
$\rho_i:G\to\operatorname{End}_K(W_i)$ be actual representations on
finite-dimensional spaces, indexed by a finite set $I$. For natural numbers
$m_i\le c_i$, set
\[
 U_m=\prod_{i\in I}\prod_{k<m_i}W_i,
 \qquad U_c=\prod_{i\in I}\prod_{k<c_i}W_i.
\]
These finite products realize the corresponding finite direct sums. The
representation on $U_m$ acts componentwise. Define $j:U_m\to U_c$ by retaining
the first $m_i$ copies and filling every remaining coordinate with zero.
Define $r:U_c\to U_m$ by restriction to the selected copies. Lean proves
\[
 rj=1_{U_m},\qquad
 j\rho_m(g)=\rho_c(g)j,\qquad
 r\rho_c(g)=\rho_m(g)r.
\]
The image of $j$ is exactly the subspace whose unused coordinates vanish.
Consequently $jr$ is an equivariant idempotent. A supplied equivariant
isomorphism $D:W\xrightarrow{\sim}U_c$ gives the split maps
$D^{-1}j$ and $rD$ inside the original carrier. The decomposition $D$ is an
input; these statements do not construct a decomposition of an unknown
arithmetic carrier.

The same formalization proves, for each $g\in G$,
\[
 \det\bigl(XI-\rho_m(g)\bigr)
 =\prod_{i\in I}\det\bigl(XI-\rho_i(g)\bigr)^{m_i},
 \qquad
 \operatorname{Tr}\bigl(\rho_m(g)\bigr)
 =\sum_{i\in I}(m_i:K)\operatorname{Tr}\bigl(\rho_i(g)\bigr).
\]
Constituents may have different dimensions. Zero multiplicities, zero-dimensional
constituents, and the empty index set are included. If at a permitted sample
$g$ the constituent polynomials are the images of $P_{i,g}\in E[X]$ under a
specified map $E\to K$, the summand polynomial is the image of
$\prod_iP_{i,g}^{m_i}$. Applying this independently at each coefficient field
preserves a supplied common polynomial family. This does not establish the
Hodge--Tate or Weil--Deligne conditions in strict compatibility.

Joint continuity of the constituent actions implies joint continuity of the
constructed action. The inclusion and retraction are continuous for the
product topologies, and triviality on a supplied subgroup is preserved.
Neither continuity nor triviality is inferred for a carrier whose corresponding
hypothesis has not been supplied. The maps transported by $D$ are proved
algebraically; their topological interpretation additionally requires the
appropriate continuity of $D$ and its inverse.

For the three-character example below, the construction produces an actual
three-dimensional summand inside the carrier of dimension seven. At the sample
$(0,1)$ its characteristic polynomial is $(X-1)^2(X+1)$. A separate checked
counterexample uses the coordinate inclusion $\mathbf Q\to\mathbf Q^2$, its
coordinate retraction, and the swap involution: the split law holds, but
compression sends the swap to zero and therefore fails to preserve its square.
Thus selecting whole constituent copies, with the proved equivariance, is
essential to this construction.

\paragraph{Algebraic monodromy data on the selected copies.}
Fix a field $K$ and a nonzero scalar $q\in K^\times$. On each constituent
$W_i$, supply an automorphism $F_i$ and a nilpotent endomorphism $N_i$ with
$F_iN_iF_i^{-1}=qN_i$. The scalar $q$ is common to this finite family.
The componentwise operators $F_m,N_m$ on $U_m$ satisfy the same relation;
$F_m$ is invertible and $N_m$ is nilpotent. The split maps constructed above
intertwine both operators. The proof first identifies the independently defined
kernel through an equivariant linear isomorphism
\[
 \ker N_m\simeq\prod_{i\in I}\prod_{k<m_i}\ker N_i.
\]
The relation, with $q\ne0$, proves that both $F_i$ and its inverse preserve
$\ker N_i$. Thus $F_i|_{\ker N_i}$ is an actual automorphism, and Lean proves
\[
 \det(XI-F_m|_{\ker N_m})
   =\prod_{i\in I}\det(XI-F_i|_{\ker N_i})^{m_i}.
\]
The corresponding kernel dimensions and traces are additive with multiplicity.
Common constituent kernel polynomials, when supplied through a coefficient-ring
map, give the mapped product polynomial for the summand. Whole-space Frobenius
polynomials are not substituted for these kernel polynomials.

The exact special chain has basis $e_0,\ldots,e_{n-1}$ and operators
\[
 Ne_0=0,\qquad Ne_j=e_{j-1}\ (j>0),\qquad Fe_j=aq^{-j}e_j,
 \qquad a,q\in K^\times.
\]
Lean proves $N^n=0$, and for $n>0$ also $N^{n-1}\ne0$ and
$\ker N=Ke_0$. Consequently its kernel polynomial is $X-a$, whereas its
whole-space polynomial is $\prod_{j=0}^{n-1}(X-aq^{-j})$.
The zero-length chain is included as a zero-dimensional space; the assertion
$\ker N=Ke_0$ is made only for positive length.

For $n>1$, replacing $N$ by zero retains exactly the same $F$ and the same
whole-space characteristic polynomial, while changing the degree of the kernel
polynomial from one to $n$. Both pairs still satisfy the algebraic monodromy
relation. This Lean-checked example shows that whole Frobenius data alone do
not determine the monodromy kernel polynomial. It is not asserted to be a
counterexample among the intended strictly compatible arithmetic systems.

These statements concern supplied algebraic pairs. No Weil-group or inertia
action, Hodge--Tate object, arithmetic comparison, or realization of the pair
from a Galois representation is constructed here. The polynomial convention is
$\det(XI-F|_{\ker N})$; no identification with an inertia-dependent local
Euler factor is asserted. Claim 9.2 remains partial at its full arithmetic scope.


\paragraph{Checked group, inertia and monodromy finite sums (Block 12).}
Let $K$ be a field, $G$ a group, and $\nu:G\to K^\times$ a character.
The new datum consists of an actual representation $r$ and nilpotent $N$ with
$r(g)N=\nu(g)Nr(g)$ for every $g$. For a supplied subgroup $H$, define
\[
 E(D,H)=\ker N\cap V^H.
\]
This is the intersection of independently defined kernel and actual invariant
subspaces. For normal $H$, the restricted whole-group representation and its
inverse automorphisms are constructed. A change from $\varphi$ to $\varphi h$,
with $h\in H$, leaves the restricted operator itself unchanged.

For a finite family of finite-dimensional constituents and natural multiplicities,
the actual local space of the constructed sum is linearly equivalent to the sum of the constituent local spaces, equivariantly
for every group element. Its characteristic polynomials multiply, its dimensions
and traces add, and maps intertwining both the group and $N$ restrict to the
local spaces. The explicit bounded-copy inclusion and retraction still split
there. Finite inertia image and an open inertia kernel are preserved.

Writing $Q_D(T)$ for the reversed local characteristic polynomial, Lean proves
\[
 Q_D(T)=\det(1-T r(\varphi)\mid E(D,H)),\qquad Q_D(0)=1,
 \qquad
 Q_{\bigoplus_iD_i^{\oplus m_i}}(T)=\prod_iQ_{D_i}(T)^{m_i}.
\]
The abstract formulas do not need characteristic zero. Normality and finite
dimension appear explicitly where used. Arithmetic identification of $G,H,\nu$
and $\varphi$ remains separate. In the geometric-Frobenius convention,
$\nu(\varphi)$ is the inverse residue cardinality; it is this unit, rather than
the residue cardinality, that enters the earlier special-chain parameter.

An explicit action of $C_2\times\mathbb Z$ has $\nu(\varphi)=1/3$,
$F=\operatorname{diag}(2,6,18)$ and $Ne_1=e_0$, $Ne_2=e_1$, $Ne_0=0$.
Trivial inertia gives local dimension one and denominator $1-2T$; scalar sign
inertia gives dimension zero and denominator one. The entire $F$ and $N$ are
equal in these two examples, and $N^2\ne0$. Thus inertia cannot be omitted.
Conversely, unramified chains of lengths one and three have equal Euler
denominators while only the first has $N=0$. Euler factors therefore do not
identify complete monodromy data.

The remaining arithmetic comparison is
$\operatorname{WD}(\bigoplus_i V_i^{\oplus m_i})\cong
\bigoplus_i\operatorname{WD}(V_i)^{\oplus m_i}$ for the actual attached functor,
intertwining the Weil action and monodromy, together with the intended coefficient
and Hodge--Tate comparisons. These are not supplied by the abstract datum.
Claim 9.2 remains partial at its full arithmetic scope.

Recovered multiplicities determine the semisimple isomorphism type relative
to the specified constituents, not a canonical embedding or projector inside
a repeated isotypic component. For example, the two projectors
$\operatorname{diag}(1,0)$ and $\operatorname{diag}(0,1)$ on a two-dimensional
trivial representation are distinct equivariant idempotents with trace one.

\begin{figure}[H]
\centering
\includegraphics[width=0.72\linewidth]{14_frobenius_separator_matrix.pdf}
\caption{A finite set of trace evaluations separates three compatible constituents and recovers their multiplicities.}
\end{figure}

\begin{numericalcertificate}[Multiplicity recovery]
The original companion matrix was not retrieved or rerun. A new exact
realization uses the trivial character and two independent sign characters
of $(\mathbf Z/2\mathbf Z)^2$, evaluated at $(0,0),(1,0),(0,1)$:
\[
 C=\begin{pmatrix}1&1&1\\1&-1&1\\1&1&-1\end{pmatrix},
 \qquad \det C=4,
 \qquad C\begin{pmatrix}2\\0\\1\end{pmatrix}
       =\begin{pmatrix}3\\3\\1\end{pmatrix}.
\]
Lean checks that these are the traces of actual one-dimensional
representations, recovers $(2,0,1)$ uniquely, rejects a target whose unique
solution is $(1/2,0,0)$, and rejects $(4,0,1)$ under capacities $(3,2,2)$.
The separator becomes singular modulo two. These are finite-image
representation checks, not fabricated Frobenius data for an arithmetic carrier.
\end{numericalcertificate}

The theorem is substantial only when a suitable compatible carrier and its irreducible compatible constituents are already known.  It does not assert that every geometric representation embeds into such a carrier.


\paragraph{Finite-exponential reconstruction and a cubic inverse (Block 13).}
Let $K$ be a characteristic-zero field and let $(r,N)$ be supplied group data
with nilpotent $N$ and $r(g)N=\nu(g)Nr(g)$. For a supplied scalar map satisfying
\[
t(gh)=t(g)+\nu(g)t(h),
\]
the formalization constructs
\[
\rho(g)=\exp(t(g)N)r(g),\qquad
\exp(aN)=\sum_{j=0}^{d-1}\frac{a^j}{j!}N^j\quad(N^d=0).
\]
Lean proves the representation laws, actual finite-copy and split-map
compatibility, and joint continuity from continuity of $N$, $r$ and $t$ with
the stated continuous vector-space operations. The monodromy fixed kernel
intersected with subgroup invariants is unchanged by this twist. For normal
inertia, its restricted action and finite-dimensional Euler denominator agree.
Twisting again by $-t$ recovers the datum with $N$ still supplied.

There is also an extracted inverse on the explicit class $N^3=0$:
\[
\log_2 U=(U-I)-\tfrac12(U-I)^2,\qquad
 a^{-1}\log_2(\exp(aN))=N\quad(a\ne0).
\]
Both exponential/logarithmic inverse identities are checked on the cubic
nilpotent and unipotent loci. If a supplied common sample $u$ has
$r_D(u)=r_E(u)=I$ and $t(u)\ne0$, an intertwiner of the reconstructed cubic
actions recovers both original group equivariance and monodromy equivariance.
The general nilpotence exponential theorem is not restricted to this cubic
class, but the extracted inverse proved here is.

The rational affine-group example has $u=(1,1)$, $d=(0,1/3)$ and
\[
\rho(u)=\begin{pmatrix}1&1&1/2\\0&1&1\\0&0&1\end{pmatrix},
\qquad \rho(d)=\operatorname{diag}(1,3,9).
\]
Reversing the factor order or dropping the quadratic term gives checked
failures. For the length-four chain, the quadratic logarithm has erroneous
$e_3$-to-$e_0$ coefficient $-1/3$, so its cubic restriction is substantive.
A zero-parameter sample also fails to recover $N$.

No arithmetic group, tame cocycle or required sample is constructed by these
statements. The attached comparison for every intended representation,
coefficient comparisons and Hodge--Tate/period data remain separate. Finite
inertia image and an open inertia kernel of $r$ are not asserted to persist
under this twist to $\rho$. Claim 9.2 remains partial.

\section{Finite spectral occupancy}

System completion produces arithmetic members.  It does not prove that a desired packet occurs in a global carrier.  In a fixed finite-dimensional semisimple carrier, however, occupancy is an exact projector question.

Let \(E\) be a field of characteristic zero and let \(T_1,\ldots,T_N\) be commuting operators on a finite-dimensional \(E\)-vector space \(A\), each diagonalizable over \(E\).  Thus the selected spectra split over the coefficient field.  Let \(a_i\in E\) be a selected target scalar for \(T_i\), possibly absent from its spectrum.  Define
\[
e_i
=
\prod_{\mu\in\Spec(T_i)\setminus\{a_i\}}
\frac{T_i-\mu I}{a_i-\mu},
\qquad
e_T=\prod_{i=1}^{N}e_i,
\]
with \(e_i=0\) if \(a_i\) is absent.

\begin{theorem}[Exact finite occupancy trace]
The operator \(e_T\) is the projector onto the joint eigenspace
\[
A_T
=
\{v:T_iv=a_iv\text{ for every }i\}.
\]
Moreover,
\[
\boxed{
\Tr(e_T)=\rank(e_T)=\dim A_T\in\Z_{\ge0}.
}
\]
The trace is identified with its unique nonnegative integer preimage under \(\Z\hookrightarrow E\).  Thus \(\Tr(e_T)>0\) refers to that integer, not to an ordering of an arbitrary field \(E\).  The finite packet is occupied exactly when this integer is positive.  In positive characteristic the field-valued trace can vanish for a nonzero projector; the rank test, not a positive field trace, must then be used.
\end{theorem}

\begin{proof}
For each operator, its finite set of actual eigenvalues is used without
multiplicity. The polynomial
\[
 p_{i,a}(X)=\prod_{\mu\in\Spec(T_i)\setminus\{a\}}
 (a-\mu)^{-1}(X-\mu)
\]
is used when $a$ occurs, and the zero polynomial otherwise. For an eigenvector
$T_i v=bv$, evaluation gives
\[
 p_{i,a}(T_i)v=\begin{cases}v&b=a,\\0&b\ne a.\end{cases}
\]
Because the eigenspaces span $A$, this proves that $e_i=p_{i,a_i}(T_i)$ maps
all of $A$ into the selected eigenspace and fixes that eigenspace pointwise.
It therefore proves the individual projector and image assertions, including
an absent target.

Polynomials in commuting operators commute. If projections $P,Q$ onto
$U,W$ commute, then $PQ$ maps into both $U$ and $W$, and fixes every vector
in $U\cap W$. Induction proves that the actual ordered product $e_T$ is a
projection onto $A_T$; commutation also makes the result independent of the
ordering. An empty list of conditions gives the identity.

For a projection onto a subspace, the direct sum of its image and kernel
identifies it with $I\oplus0$. Thus its trace in $E$ is the natural dimension
of its image cast into $E$. In characteristic zero this cast is injective,
so there is a unique natural integer represented by the trace. That integer
is positive exactly when the joint eigenspace contains a nonzero vector.
\end{proof}

\paragraph{Naturality and finite hypotheses.}
If $F T_i=U_i F$ and the source operators split, then
$F e_{T,a}=e_{U,a}F$. This follows first on source eigenvectors and then on
their span. In particular, a change of basis conjugates the constructed
projector and preserves its trace. Adding conditions decreases the actual
joint eigenspace and cannot increase its dimension. These statements do not
assume an orthonormal basis or a positive definite inner product.

\paragraph{A specified finite realization.}
On $\Q^4$, take the diagonal operators with diagonal lists
\[
(1,1,1,0),\quad(2,2,0,0),\quad(3,0,0,0),\quad(0,4,4,4)
\]
and target scalars $(1,2,3,4)$. All four targets occur individually. Imposing
these conditions successively gives actual projector ranks $3,2,1,0$.
The last packet is empty despite individual feasibility. The formal code
checks these operators, projector formulas, traces and dimensions. This is
a newly specified realization, not a claim to have recovered the original
companion implementation.

\paragraph{Why the hypotheses cannot be omitted.}
For
\[
 P=\begin{pmatrix}1&0\\0&0\end{pmatrix},\qquad
 Q=\tfrac12\begin{pmatrix}1&1\\1&1\end{pmatrix},
\]
each operator is idempotent, but $PQ$ is not idempotent and has trace $1/2$.
For the Jordan block $J=\left(\begin{smallmatrix}1&1\\0&1\end{smallmatrix}\right)$,
the singleton-spectrum Lagrange polynomial is $1$, so the formula produces
the identity although the eigenspace at $1$ is only a line. Both failures
are exact formal counterexamples to dropping the corresponding hypotheses,
not counterexamples to the theorem with those hypotheses retained.


\begin{figure}[H]
\centering
\includegraphics[width=0.77\linewidth]{15_finite_occupancy_ranks.pdf}
\caption{A consistent sequence of finite conditions narrows the carrier from ranks \(3\) to \(2\) to \(1\).  One inconsistent condition forces the joint projector to zero.}
\end{figure}

\begin{conjecture}[Universal occupancy]
For every admissible finite arithmetic packet in the intended global class, there exists a finite-dimensional carrier in which the associated joint projector has positive trace.
\end{conjecture}

This conjecture is not a consequence of the projector formula.  The formula measures occupancy exactly once the carrier is fixed.  It cannot force the trace to be positive.

\section{A conditional system-completion criterion}

The preceding results can be assembled without promoting their missing hypotheses to theorems.

\begin{theorem}[Residual-system completion criterion]
Fix an admissible arithmetic class.  Suppose:
\begin{enumerate}
\item every object has a unique residual root in an exhaustive residual moduli space;
\item a common algebraic spine is defined for the desired system;
\item every coefficient-prime residual companion exists;
\item every companion forest has nonempty finite levels, finite branching, and inverse-limit closed local conditions;
\item the required horizontal transports preserve the common spine and the selected local deformation conditions.
\end{enumerate}
Then the desired strictly compatible arithmetic system exists.  If, in addition, the system is a summand of a known compatible carrier, the summand is recovered at every coefficient prime by finitely many Frobenius traces.
\end{theorem}

\begin{proof}
The first four hypotheses give the system by adelic finite-depth completion.  The transport hypothesis identifies any roots that are meant to be generated horizontally without changing the prescribed spine.  The final assertion is the compatible-carrier summand theorem.
\end{proof}

\begin{corollary}[Exact location of the remaining existence problem]
Inverse-limit reconstruction alone is not a nonemptiness theorem.  The existence requirements must retain both residual companions and liftability through every finite depth:
\[
\boxed{
\begin{aligned}
\text{companion existence: }&X_{\lambda,1}(\Compat)\ne\varnothing
\text{ for every }\lambda;\\
\text{finite-depth liftability: }&X_{\lambda,m}(\Compat)\ne\varnothing
\text{ for every }\lambda,m\ge2;\\
\text{finite occupancy: }&\Tr(e_T)>0
\text{ for every required global packet.}
\end{aligned}}
\]
Formal reconstruction does not create a missing residual root, repair an empty higher level, or force global occupancy.  Nonempty finite levels with finite branching supply a branch, after which inverse-limit closure reconstructs an admissible object.
\end{corollary}


\paragraph{An occupied node can have no child.}
The actual order-two representation with generator
\[
 A=\begin{pmatrix}1&2\\2&3\end{pmatrix}\pmod4
\]
satisfies $A^2=I$ and $\det(XI-A)=X^2-1$. It has no lift modulo $8$ with
the same characteristic polynomial. Indeed, every lift has entries congruent
to $1,2,2,3$ modulo $4$. Trace zero then forces its determinant to be $3$
modulo $8$, whereas $X^2-1$ requires determinant $-1=7$ modulo $8$.
Both the occupied node and the absence of any next-level matrix or representation
with the prescribed polynomial are Lean-checked. Other branches of the same
unrestricted involution forest survive, so this example does not contradict
finite-branch compactness.

\paragraph{Local properties need their own limit argument.}
Over a domain, for a nonzero $\pi\in I$, the representations of the additive
integer group given by
\[
 \rho_n(g)=\begin{pmatrix}1&\pi^n g\\0&1\end{pmatrix}
\]
are all nontrivial and converge pointwise in the ideal-adic topology to the
trivial representation. The convergence and the failure of limit preservation
for nontriviality are formally checked. Thus not every property of approximating
integral representations can be asserted for a limit. This example is not
claimed to disprove stability for any particular intended Galois local condition.

For the completed congruence problem, the eleventh block proves the full
conditional equivalence with residual nonemptiness, nonemptiness at every higher
depth, and nonzero trace for \emph{every} required supplied packet. The packet
index may be infinite; each individual list of split commuting conditions is
finite, over a characteristic-zero field. This is an equivalence with actual
realizing continuous representations and the supplied local conditions, not a
positive-occupancy theorem. Claim 11.2 remains partial only at the full arithmetic
interpretation: the required arithmetic objects and packet carriers must be
identified with the formal data. No construction of an unknown carrier is
asserted by this conditional boundary statement.

\begin{conjecture}[Natural finite seed rank]
For the admissible residual transport category attached to a fixed arithmetic class, the directed seed rank is finite.  A one-seed theorem additionally requires a source component from which every required object is reachable and transports preserving deformation and occupancy data.  For a finite condensation DAG, having exactly one source implies that reachability condition.  For an infinite category, the existence of a unique source alone does not establish it.
\end{conjecture}

\section{Worked finite models}

\subsection{Smooth branch}

For \(q=3\) and formal dimension \(d=2\), the finite levels have cardinalities
\[
1,9,81,729,6561,\ldots.
\]
The seed zeta function is
\[
Z(T)=\frac{T}{1-9T}.
\]

\subsection{Weighted dyadic chart}

For weights \((1,2,4)\), the counts are
\[
1,2,8,32,256,2048,\ldots.
\]
The first coordinate is active immediately, the second one level later, and the third at depth five.  The exact rational function generated by the code is
\[
Z_{(1,2,4)}(T)
=
\frac{T(1-6T-8T^2-32T^3)}{1-8T}.
\]
The negative signs in this presentation arise from moving the denominator to \(1-8T\); every power-series coefficient is positive.

\subsection{Singular branch}

For \(x^2=0\) over a residue field of size three, the counts are
\[
1,3,3,9,9,27,27,81,\ldots,
\]
and
\[
Z(T)=\frac{T+3T^2}{1-3T^2}.
\]
The exact branching dimension is \(1/2\).

\subsection{Directed transport}

The finite graph in \cref{fig:transport-graph} has three source components.  The source-component theorem and brute-force enumeration both return the minimum seed count three.

\begin{figure}[H]
\centering
\includegraphics[width=0.84\linewidth]{16_proof_status.pdf}
\caption{The finite and formal statements proved in this paper are separated from the two remaining universal existence assertions.}
\label{fig:proof-status}
\end{figure}

\section{Computational certificates}

The code uses exact integer, rational, symbolic, and modular arithmetic for all mathematical identities.  Floating-point arithmetic is used only to render the figures and to report a regression estimate alongside an exact branching dimension.

\begin{table}[H]
\centering
\caption{Selected exact certificates.}
\begin{tabular}{lr}
\toprule
Certificate & Result\\
\midrule
Residual packet count in the mod-three model & 3\\
Compatible base-five inverse-limit branch & yes\\
First dyadic sign separation depth & 2\\
Source strongly connected components & 3\\
Brute-force minimum directed seed count & 3\\
Cohomology packet reconstruction residual & 0\\
Compatible common-spine test & passed\\
Inconsistent spine detected & yes\\
Recovered carrier multiplicities & \((2,0,1)\)\\
Finite occupancy ranks & \((3,2,1,0)\)\\
\bottomrule
\end{tabular}
\end{table}

The complete automated suite contains \(28\) tests.  They cover residual partition, no crossing, inverse-limit digits, weighted and singular zeta coefficients, directed source components, cohomology packets, companion forests, common-spine consistency, finite Frobenius separation, multiplicity recovery, and exact occupancy projectors.

\section{Claim ledger}

\begin{longtable}{p{0.34\linewidth}p{0.18\linewidth}p{0.40\linewidth}}
\toprule
Statement & Status & Scope or boundary\\
\midrule
\endfirsthead
\toprule
Statement & Status & Scope or boundary\\
\midrule
\endhead
Residual partition & proved & any moduli functor with well-defined reduction\\
No crossing & proved & operations trivial at residual depth\\
Inverse-limit reconstruction & proved & represented complete deformation problem\\
Finite-branch compactness & proved & finite nonempty levels and finite branching\\
Smooth and weighted seed zeta formulas & proved & stated formal coordinate models\\
Singular power zeta and dimension & proved & scalar equation \(x^e=0\)\\
Cohomology packet decomposition & proved & chain projectors commuting with the differential\\
Directed minimum-seed theorem & proved & finite directed transport graph\\
Adelic finite-depth completion & proved & common spine, all finite levels occupied, closed local conditions\\
Compatible-carrier summand completion & proved & known semisimple compatible carrier and finite separation\\
Finite occupancy trace & Lean checked & supplied commuting split operators; characteristic-zero trace criterion\\
Universal residual companions & not proved & requires nonempty root at every coefficient prime\\
Universal compatible-system completion & not proved & algebraic spine and companion occupancy remain necessary\\
Universal automorphic occupancy & not proved & requires positive global packet traces\\
Unrestricted reciprocity & not claimed & belongs to later local and global papers\\
\bottomrule
\end{longtable}

\section{Conclusion}

The residual seed theory becomes precise once vertical and horizontal generation are separated.  Positive-depth perturbations can be complete within one formal component.  They cannot change the residual root.  The universal arithmetic seed is therefore a forest rather than one branch:
\[
\boxed{
\text{residual roots}
\longrightarrow
\text{finite deformation levels}
\longrightarrow
\text{inverse-limit branches}.
}
\]

Projectors add three useful structures to this forest.  They split tangent and obstruction complexes into exact subproblems, they measure delayed dyadic visibility through weights, and they turn finite carrier occupancy into an integer trace.  Directed transport adds a fourth: it converts the informal question ``how many seeds are needed?'' into the source-component count of a condensation graph.

The system-level boundary is now narrow enough to state without ambiguity.  One must construct a common algebraic spine, populate every residual companion forest, lift those companions through all depths, and prove positive occupancy of the required global packets.  The present paper proves the formal reconstruction, counting, transport, compactness, carrier, and finite-projector theorems used in that reduction.  It does not assert the missing universal existence statements.

\appendix

\section{Algorithms implemented in the companion package}

\subsection{Residual towers}

The residual module partitions coordinates modulo \(p\), verifies compatibility
\[
x_{m+1}\equiv x_m\pmod{p^m},
\]
and reconstructs the digit \(a_{m-1}\) by
\[
a_{m-1}=\frac{x_m}{p^{m-1}}\pmod p.
\]

\subsection{Transport condensation}

The transport module computes strongly connected components by two depth-first passes, forms the condensation graph, finds all source components, and selects one representative from each.  A brute-force subset search is included only as an independent finite certificate.

\subsection{Companion forests}

For an integral polynomial \(f\), the finite root forest is obtained by exact enumeration of
\[
f(x)\equiv0\pmod{p^m}
\]
and by connecting a child \(x_{m+1}\) to its reduction modulo \(p^m\).  This is intended for small exact examples, not large deformation rings.

\subsection{Cohomology packets}

Modular Gaussian elimination computes the ranks of the differentials.  For
\[
C^{i-1}\xrightarrow{d_{i-1}}C^i\xrightarrow{d_i}C^{i+1},
\]
the finite dimension is evaluated as
\[
\dim H^i
=
\dim C^i-\rank d_i-\rank d_{i-1}.
\]
Packet subcomplexes are computed separately and checked against the total.

\subsection{Finite occupancy}

For each commuting semisimple operator, a Lagrange polynomial constructs the exact spectral idempotent.  Their product gives the joint packet projector.  Symbolic rational arithmetic verifies idempotence and the equality of trace and rank.

\section{Reproducibility boundary}

The companion code is a finite certificate library.  It does not construct universal Galois deformation rings, compute general continuous group cohomology, prove representability, discover compatible systems, or evaluate automorphic trace formulas.  Those are mathematical inputs or later existence problems.  The release is designed to make every finite example and every displayed count independently reproducible without confusing computation with proof.


\section*{Current formalization boundary}
The smooth formal-model classification, its finite point count, its integral digit streams, and its rational seed zeta identity are now checked in Lean. The actual representing-ring algebra isomorphism remains a supplied hypothesis; neither smoothness nor Galois representability is inferred from numerical counts. All local maps are included in the classification, and locality and actual maximal-ideal-adic continuity are separately proved. Arbitrary choices of parameter representatives do not add phantom maps.

The cumulative ledger contains fifteen checked source proof claims and five partial claims, with no untouched proof claims. Exact finite occupancy is checked. Claims 8.2, 9.1, 9.2, 11.1 and 11.2 retain the arithmetic gaps recorded in the current ledger. Finite-exponential reconstruction and cubic logarithmic recovery advance the local comparison without proving arithmetic attachment or universal positive occupancy. Universal occupancy and natural finite seed rank remain conjectures. The original companion package has not been rerun. The release provides separate exact checks and a cumulative build of all project production modules. No new PDF was typeset.


\paragraph{Finite-group comparison with precision loss (Block 17).}
Let $R$ be a commutative domain, $\Gamma$ a finite group, and $\rho,\tau$
actual integral matrix representations. Suppose $a=|\Gamma|\ne0$ in $R$ and
$a\mid\pi^e$. An equivariant matrix modulo $(\pi)^{e+k}$ lifts to an
integral equivariant matrix with the prescribed reduction modulo $(\pi)^k$.
Indeed, lift its entries to $A$ and form
\[
S=\sum_{g\in\Gamma}\rho(g)A\tau(g^{-1}).
\]
The finite equations imply $S-aA=\pi^{e+k}E$ with $E$ integral. Choose $b$
with $\pi^e=ab$ and put $P=A+b\pi^kE$. Then $aP=S$, so cancellation of
$a\ne0$ proves equivariance. Over a local ring, if $\pi$ is a nonunit,
$k>0$ and the supplied comparison is invertible, $P$ is an integral unit.
Neither completeness nor a finite residue field is needed for this step.

For a DVR, $e=v_\pi(|\Gamma|)$ is derived from the actual group order.
The single precision $e+1$, independent of rank, detects integral
isomorphism. It makes reduction injective on integral isomorphism classes,
and finite coefficient quotients give finitely many integral classes in
fixed rank. This is a formalization of classical finite-group rigidity
(Maranda, 1953, Theorem~2), with an explicit proof of the comparison's
precision loss.

For complete local coefficients with finite quotients, downward-compatible
allowed representations occupied at every precision, and the required
finite-group type at depth $e+1$, reconstruction followed by this comparison
produces the required integral representation. Continuity is proved for a
finite discrete source and for pullback through a fixed continuous map to
that source. Occupancy at only one precision is insufficient. The result
lifts comparisons between actual integral representations; it does not lift
arbitrary quotient representations or supply the missing arithmetic packet
interpretation.

For $C_2$ over $\mathbb Z_2$, take
\[
\tau(g)=\begin{pmatrix}1&0\\0&-1\end{pmatrix},\qquad
\rho(g)=\begin{pmatrix}1&2\\0&-1\end{pmatrix},\qquad
A=\begin{pmatrix}1&1\\0&1\end{pmatrix}.
\]
The matrix $A$ intertwines modulo $4$ and its average is
$P=\left(\begin{smallmatrix}1&-1\\0&1\end{smallmatrix}\right)$.
No exact intertwiner can retain $A$ modulo $4$, although $P$ retains it
modulo $2$. The actual trivial and sign representations also agree modulo
$2$ but are not integrally isomorphic. These checked examples show why
residual-only detection and full preservation of the input precision cannot
be asserted here.


\paragraph{Lifting a summand inside a supplied finite-group carrier (Block 18).}
Let $R$ be a commutative integral domain complete and separated for the
$(\pi)$-adic filtration. Let $\Gamma$ be a finite group, and let
$\rho:\Gamma\to\mathrm{GL}_n(R)$ be an actual integral representation.
Assume that $a=|\Gamma|\ne0$ in $R$ and $a\mid\pi^e$. For $k>0$, suppose
\[
\bar E\in M_n\bigl(R/(\pi)^{e+k}\bigr),\qquad
\bar E^2=\bar E,\qquad
\bar\rho(g)\bar E=\bar E\bar\rho(g)\quad(g\in\Gamma).
\]
There is an actual integral matrix $E$ satisfying
\[
E^2=E,\qquad \rho(g)E=E\rho(g),\qquad
E\bmod(\pi)^k=\bar E\bmod(\pi)^k.
\]
No integral summand or coherent family of quotient projectors is assumed.
No local-ring assumption or finiteness of the residue field is needed.

To prove this, apply the preceding averaging correction with both
representations equal to $\rho$. This produces an exact equivariant matrix
$A$ with the prescribed reduction at depth $k$, so that
$A^2-A\in(\pi)^kM_n(R)$. Averaging alone need not preserve idempotence.
Use the integral polynomial
\[
f(A)=3A^2-2A^3.
\]
The identities
\[
f(A)-A=(1-2A)(A^2-A),\qquad
f(A)^2-f(A)=(A^2-A)^2(4A^2-4A-3)
\]
show that successive corrections preserve the retained reduction and improve
the idempotence defect. In these formulas scalar constants denote scalar
matrices. The Lean proof establishes at least one additional ideal digit at
each step, constructs the resulting compatible family of finite-level matrices, and uses
completeness and separatedness to recover $E$. Every matrix commuting with
$A$ commutes with $E$. There is no division by $2$, so the correction also
works in residue characteristic $2$.

The actual submodules $U=\operatorname{im}E$ and $W=\ker E$ are invariant,
and addition gives an equivariant isomorphism
\[
U\oplus W\simeq R^n.
\]
The inclusion $i:U\hookrightarrow R^n$ and retraction $r(v)=Ev$ are
constructed as linear maps and satisfy $ri=1_U$ and $ir=E$. Their
equivariance is proved. Coordinatewise reduction maps $U$ onto
$\operatorname{im}(\bar E\bmod(\pi)^k)$ and maps $W$ onto
$\ker(\bar E\bmod(\pi)^k)$. Both equalities are surjectivity statements
onto the actual submodules over the finite-level quotient, not merely containment. The same
projector defines a summand after pullback through any group homomorphism
to $\Gamma$.

For $R=\mathbb Z_2$ and the sign carrier
$\rho(s)=\operatorname{diag}(1,-1)$, the projector
\[
\bar E=\begin{pmatrix}1&2\\0&0\end{pmatrix}\pmod4
\]
is equivariant and the theorem constructs a nonzero proper equivariant
projector retaining its reduction modulo $2$. No exact equivariant
endomorphism retains this prescribed matrix modulo $4$. At the insufficient
input depth $1$, the projector
$\left(\begin{smallmatrix}1&1\\0&0\end{smallmatrix}\right)\pmod2$
has no exact equivariant lift retaining its reduction. These are statements
about the supplied carrier and prescribed projectors, not about nonexistence
of all integral summands.

This block constructs one summand and its complement inside a supplied
finite-group carrier. It advances the carrier step in Claim 9.2. It does
not identify the intended arithmetic packets, construct the arithmetic
Weil--Deligne attachment, establish the Hodge--Tate comparisons, or prove
universal occupancy. The full arithmetic ledger remains fifteen checked
claims and five partially formalized claims.


\paragraph{Simultaneous projector-family lifting and rigidity (Block 19).}
Keep the complete domain, supplied finite-group carrier and group-order
hypotheses of Block 18. Let $J$ be any finite label set and $k>0$. Suppose
$(\bar E_i)_{i\in J}$ is a complete orthogonal family of equivariant
idempotents modulo $(\pi)^{e+k}$. There exists an integral equivariant
family $(E_i)_{i\in J}$ satisfying
\[
E_i^2=E_i,\qquad E_iE_j=0\ (i\ne j),\qquad
\sum_{i\in J}E_i=1,\qquad
E_i\bmod(\pi)^k=\bar E_i\bmod(\pi)^k.
\]
The loss is at most $e$, independent of $|J|$ and the ambient matrix rank.
There is no local-ring or finite-residue-field hypothesis for existence.
The formal theorem also treats empty label sets whenever their input
equations can hold.

Average every supplied matrix once. Inductively let $C$ be the complement
of the previously lifted projectors and correct the next averaged matrix
inside the corner defined by $C$. The integral polynomial correction
preserves the retained reduction and commutation with $C$ and the carrier.
Multiplication by $C$ gives the required supported projector. This consumes
no further precision. Once every label has been lifted, the complementary
idempotent reduces to zero at positive depth. Adic separatedness forces it
to vanish, proving the sum-to-identity equation.

The actual invariant images $U_i=\operatorname{im}E_i$ give an equivariant
linear isomorphism
\[
\bigoplus_{i\in J}U_i\simeq R^n,
\qquad (v_i)_i\longmapsto\sum_i v_i,
\qquad v\longmapsto(E_iv)_i.
\]
Lean represents the finite direct sum by the finite product of the image
submodules. The existing coordinatewise reduction theorem maps every
$U_i$ onto the corresponding reduced projector image. No chosen bases or
unproved freeness of the images are used.

For rigidity alone, let $R$ be a commutative local ring, let $\pi$ be a
nonunit, and let $E,F$ be complete equivariant families on the same carrier
with $E_i\equiv F_i\pmod{(\pi)^k}$ for every label and $k>0$. Then
\[
U=\sum_{i\in J}F_iE_i,\qquad
U\equiv1\pmod{(\pi)^k},\qquad UE_i=F_iU.
\]
Reduction reflects invertibility in this setting, so $U$ is an integral
unit and $UE_iU^{-1}=F_i$. It commutes with the carrier and induces actual
equivariant isomorphisms $\operatorname{im}E_i\simeq\operatorname{im}F_i$.
This rigidity theorem requires neither a finite acting group, a domain,
nor completeness. More generally a unit-reflecting coefficient map suffices.
The combined result proves existence and uniqueness up to this labeled
conjugacy; exact uniqueness of the matrices is false.

For the continuous $\mathbb Z_2$ carrier
$\rho(s)=\operatorname{diag}(1,-1,1)$, conjugate the three coordinate
projectors by $1+2e_{01}$ and reduce modulo $4$. They form a complete
equivariant family. The general theorem constructs three nonzero integral
summands retaining this family modulo $2$. No exact equivariant family
retains the whole prescribed family modulo $4$. Separate checked examples
show that independent equivariant idempotent lifts can lose orthogonality,
and that distinct complete families with the same reduction can be related
by the asserted identity-congruent conjugacy.

Block 19 completes this finite-group projector-family step. The actual
arithmetic packets, Galois-to-Weil--Deligne attachment, coefficient and
Hodge--Tate/period comparisons, and arithmetic Frobenius coverage remain
outside this result. The full arithmetic ledger remains fifteen checked
claims and five partial claims; universal occupancy remains conjectural.

\end{document}
